% tikz-setup.tex -- einmal in der Präambel: % tikz-setup.tex -- einmal in der Präambel: % tikz-setup.tex -- einmal in der Präambel: % tikz-setup.tex -- einmal in der Präambel: \input{tikz-setup}
\usepackage{tikz}
\usetikzlibrary{arrows.meta, positioning, calc, fit, shapes.geometric, decorations.pathreplacing}

\definecolor{figblue}{HTML}{3B6EA5}
\definecolor{figorange}{HTML}{D9822B}
\definecolor{figgreen}{HTML}{4C9A6A}
\definecolor{figpurple}{HTML}{7B5EA7}
\definecolor{figgray}{HTML}{808080}

\tikzset{
  figbox/.style={draw=black!70, rounded corners=2pt, fill=white, align=center,
    inner sep=4pt, font=\small, minimum height=9mm},
  sblue/.style={fill=figblue!15, draw=figblue},
  sorange/.style={fill=figorange!15, draw=figorange},
  sgreen/.style={fill=figgreen!15, draw=figgreen},
  spurple/.style={fill=figpurple!15, draw=figpurple},
  figarrow/.style={-{Stealth[length=2.2mm]}, line width=0.8pt, black!75},
  figlabel/.style={font=\footnotesize, text=black!75, align=center},
  figneuron/.style={circle, draw=black!70, fill=white, minimum size=5mm, inner sep=0pt},
  imgcoords/.style={shift={(#1.south west)},
    x={($(#1.south east)-(#1.south west)$)},
    y={($(#1.north west)-(#1.south west)$)}},
}

% Platzhalter/echtes Münzbild (Dateiname anpassen); in der Preview vorher definiert
\providecommand{\coinpic}[1]{\includegraphics[width=#1]{coin3545_obverse.jpeg}}

\usepackage{tikz}
\usetikzlibrary{arrows.meta, positioning, calc, fit, shapes.geometric, decorations.pathreplacing}

\definecolor{figblue}{HTML}{3B6EA5}
\definecolor{figorange}{HTML}{D9822B}
\definecolor{figgreen}{HTML}{4C9A6A}
\definecolor{figpurple}{HTML}{7B5EA7}
\definecolor{figgray}{HTML}{808080}

\tikzset{
  figbox/.style={draw=black!70, rounded corners=2pt, fill=white, align=center,
    inner sep=4pt, font=\small, minimum height=9mm},
  sblue/.style={fill=figblue!15, draw=figblue},
  sorange/.style={fill=figorange!15, draw=figorange},
  sgreen/.style={fill=figgreen!15, draw=figgreen},
  spurple/.style={fill=figpurple!15, draw=figpurple},
  figarrow/.style={-{Stealth[length=2.2mm]}, line width=0.8pt, black!75},
  figlabel/.style={font=\footnotesize, text=black!75, align=center},
  figneuron/.style={circle, draw=black!70, fill=white, minimum size=5mm, inner sep=0pt},
  imgcoords/.style={shift={(#1.south west)},
    x={($(#1.south east)-(#1.south west)$)},
    y={($(#1.north west)-(#1.south west)$)}},
}

% Platzhalter/echtes Münzbild (Dateiname anpassen); in der Preview vorher definiert
\providecommand{\coinpic}[1]{\includegraphics[width=#1]{coin3545_obverse.jpeg}}

\usepackage{tikz}
\usetikzlibrary{arrows.meta, positioning, calc, fit, shapes.geometric, decorations.pathreplacing}

\definecolor{figblue}{HTML}{3B6EA5}
\definecolor{figorange}{HTML}{D9822B}
\definecolor{figgreen}{HTML}{4C9A6A}
\definecolor{figpurple}{HTML}{7B5EA7}
\definecolor{figgray}{HTML}{808080}

\tikzset{
  figbox/.style={draw=black!70, rounded corners=2pt, fill=white, align=center,
    inner sep=4pt, font=\small, minimum height=9mm},
  sblue/.style={fill=figblue!15, draw=figblue},
  sorange/.style={fill=figorange!15, draw=figorange},
  sgreen/.style={fill=figgreen!15, draw=figgreen},
  spurple/.style={fill=figpurple!15, draw=figpurple},
  figarrow/.style={-{Stealth[length=2.2mm]}, line width=0.8pt, black!75},
  figlabel/.style={font=\footnotesize, text=black!75, align=center},
  figneuron/.style={circle, draw=black!70, fill=white, minimum size=5mm, inner sep=0pt},
  imgcoords/.style={shift={(#1.south west)},
    x={($(#1.south east)-(#1.south west)$)},
    y={($(#1.north west)-(#1.south west)$)}},
}

% Platzhalter/echtes Münzbild (Dateiname anpassen); in der Preview vorher definiert
\providecommand{\coinpic}[1]{\includegraphics[width=#1]{coin3545_obverse.jpeg}}

\usepackage{tikz}
\usetikzlibrary{arrows.meta, positioning, calc, fit, shapes.geometric, decorations.pathreplacing}

\definecolor{figblue}{HTML}{3B6EA5}
\definecolor{figorange}{HTML}{D9822B}
\definecolor{figgreen}{HTML}{4C9A6A}
\definecolor{figpurple}{HTML}{7B5EA7}
\definecolor{figgray}{HTML}{808080}

\tikzset{
  figbox/.style={draw=black!70, rounded corners=2pt, fill=white, align=center,
    inner sep=4pt, font=\small, minimum height=9mm},
  sblue/.style={fill=figblue!15, draw=figblue},
  sorange/.style={fill=figorange!15, draw=figorange},
  sgreen/.style={fill=figgreen!15, draw=figgreen},
  spurple/.style={fill=figpurple!15, draw=figpurple},
  figarrow/.style={-{Stealth[length=2.2mm]}, line width=0.8pt, black!75},
  figlabel/.style={font=\footnotesize, text=black!75, align=center},
  figneuron/.style={circle, draw=black!70, fill=white, minimum size=5mm, inner sep=0pt},
  imgcoords/.style={shift={(#1.south west)},
    x={($(#1.south east)-(#1.south west)$)},
    y={($(#1.north west)-(#1.south west)$)}},
}

% Platzhalter/echtes Münzbild (Dateiname anpassen); in der Preview vorher definiert
\providecommand{\coinpic}[1]{\includegraphics[width=#1]{coin3545_obverse.jpeg}}
